\documentclass[a4paper,12pt]{article}
\usepackage[utf8]{inputenc}
\usepackage[T1]{fontenc}
\usepackage[ngerman]{babel}
\usepackage{amsmath}
\usepackage{amssymb}
\usepackage{enumitem}
\usepackage{geometry}
\geometry{a4paper, margin=2.5cm}

\title{Einsendeaufgaben -- Lektion 6\\[0.5em]
\large Modul 61111: Mathematische Grundlagen}
\author{}
\date{}

\begin{document}
\maketitle

\section*{Aufgabe 6.1}

Sei $x \geq 0$ beliebig. Wir definieren
\[
  f\colon [0, x] \to \mathbb{R}, \quad
  t \mapsto \ln(1+t) - \frac{t}{\sqrt{1+t}}.
\]

Die Ableitung von $f$ ist:
\[
  f'(t) = \frac{1}{1+t} - \frac{\sqrt{1+t} + t \cdot \dfrac{1}{2\sqrt{1+t}}}{1+t}
  = \frac{1 - \sqrt{1+t} - t \cdot \dfrac{1}{\sqrt{1+t}}}{1+t}.
\]

Für alle $t \geq 0$ gilt $f'(t) \leq 0$, da der Nenner von $f'(t)$ immer positiv ist
und der Zähler wegen $\sqrt{1+t} \geq 1$ immer negativ ist.

Da $x > 0$, existiert nach dem Mittelwertsatz ein $t \in (0, x)$ mit
\[
  f'(t) = \frac{f(x) - f(0)}{x - 0} = \frac{f(x)}{x}.
\]

Mit $f'(t) \leq 0$ folgt:
\[
  0 \geq f'(t) = \frac{f(x)}{x}.
\]

Da $x > 0$ gilt, folgt:
\[
  0 > f(x) = \ln(1+x) - \frac{x}{\sqrt{1+x}}.
\]

Nach Umstellung gilt:
\[
  \ln(1+x) < \frac{x}{\sqrt{1+x}}. \qquad \square
\]

\end{document}
